% 数理逻辑（综述）
% license CCBYSA3
% type Wiki

本文根据 CC-BY-SA 协议转载翻译自维基百科\href{https://en.wikipedia.org/wiki/Mathematical_logic}{相关文章}。

数学逻辑（Mathematical Logic）是数学中对形式逻辑的研究。其主要分支包括模型论、证明论、集合论与递归论（又称可计算性理论）。数学逻辑的研究通常探讨逻辑形式系统的数学性质，例如其表达能力与推理能力。然而，它也可以用于刻画数学推理的正确性，或为数学建立理论基础。

自其诞生以来，数学逻辑既推动了数学基础研究的发展，也受到其驱动。该领域的起源可追溯至19世纪晚期，当时几何、算术与分析的公理化体系相继建立。20世纪初，大卫·希尔伯特（David Hilbert）提出了著名的“希尔伯特计划”（Hilbert’s Program），旨在证明数学基础理论的一致性。库尔特·哥德尔（Kurt Gödel）、格哈德·根岑（Gerhard Gentzen）等人的研究部分地解决了这一计划，同时也揭示了证明一致性问题的复杂性。

集合论的研究表明，几乎所有通常的数学内容都可以用集合语言形式化，尽管仍存在一些定理无法在常见集合论公理系统内被证明。现代数学基础研究的重点，往往不再是寻找能涵盖全部数学的理论，而是致力于分析哪些数学部分能够在特定的形式系统中被形式化（如“反向数学”所做的那样）。
